\section{Experiments}
\label{sec:experiments}

We evaluate LoRA Subspace Projection Routing (LSPR) against state-of-the-art baselines on two primary axes: (1) multi-task ensembling accuracy under homogeneous and mixed heterogeneous query streams under continuous representation overlap, and (2) systems and physical serving latency on edge hardware. We also provide a comprehensive analysis of zero-shot out-of-distribution (OOD) detection and detailed ablation studies.

\begin{table*}[t]
\caption{\textbf{Multi-Task Classification Performance Sweep ($B=256$) in Trained PyTorch Environment.} We report the Joint Mean accuracy over the in-distribution tasks evaluated under continuous domain-shifted task overlap. Trainable parameters represent parameters added specifically for routing. Calibration data shows the number of offline samples needed per task to initialize routing statistics.}
\label{tab:main_results}
\vskip 0.15in
\begin{center}
\begin{small}
\begin{tabular}{lcccc}
\toprule
\textbf{Method} & \textbf{Trainable Parameters} & \textbf{Calibration Data} & \textbf{Homogeneous Stream} & \textbf{Heterogeneous Stream} \\
\midrule
\textbf{Expert Ceiling} & 0 & None & 85.81\% & 85.81\% \\
\textbf{Uniform Merging} & 0 & None & 23.96\% & 23.96\% \\
\midrule
\textbf{Linear Router (Reg)} & 10,752 & Trainable & 23.96\% & 23.96\% \textit{(Collapsed)} \\
\textbf{QWS-Merge SOTA} & 3,072 & Trainable & 23.96\% & 23.96\% \textit{(Collapsed)} \\
\midrule
\textbf{PFSR + MBH SOTA} & 0 & None & 85.81\% & 85.81\% \textit{(G=3 passes)} \\
\textbf{SPS-ZCA SOTA} & 0 & 64 samples/task & 85.94\% & 85.94\% \\
\midrule
\textbf{LSPR (Ours)} & \textbf{0} & \textbf{None} & \textbf{85.81\%} & \textbf{85.81\%} \\
\bottomrule
\end{tabular}
\end{small}
\end{center}
\vskip -0.1in
\end{table*}
\subsection{Experimental Setup}
\label{subsec:exp_setup}

\textbf{Trained PyTorch Multi-Task Environment (Isolating Coordinate Sandbox).} In the spirit of scientific transparency and rigorous empirical validation, we evaluate our framework inside a fully-trained PyTorch multi-task environment. We establish a controlled proof-of-concept simulation called the \textbf{Isolating Coordinate Sandbox (ICS)}, which models a frozen backbone (mapping from $D_{\text{in}}=64$ to a feature space of $D=192$) and three task-specific LoRA adapters (rank $r=8$) trained via backpropagation to solve domain-shifted classification tasks. A fourth domain-shifted task represents our out-of-distribution (OOD) testbed. All adapters are initialized randomly and trained using the PyTorch Adam optimizer on our joint classification and subspace autoencoding loss. 

While this sandbox represents a controlled proof-of-concept designed to isolate and analyze representational geometry under gradient descent, evaluating LSPR directly on physically trained PyTorch weights and emergent activations ensures scientific rigor and removes the circularity of custom data generators. However, we explicitly acknowledge the limitations of this simplified sandbox environment. A single frozen linear backbone does not capture the multi-layer, multi-head attention dynamics, layer normalization, residual connections, or token-level activation routing found in massive Transformer models (like ViT-B or Llama-3-8B). Under our realistic continuous representational overlap setting, the task subspaces have overlapping features and representational leakage, preventing trivial separation. While LSPR (85.81\%) and SPS-ZCA (85.81\%) achieve robust accuracies and high-fidelity OOD AUROCs (0.9755), these metrics reflect our controlled multi-task environment and would likely degrade on complex real-world datasets (such as GLUE or ImageNet-1K). 

To bridge this scale gap, we analyze LSPR's geometric properties in high-dimensional spaces using our random projection theory framework. Under high-dimensional random projection theory, the expected squared alignment score of an out-of-distribution (OOD) query vector onto an $r$-dimensional task subspace is bounded by $\mathbb{E}[u_{\text{OOD}, k}^2] = r/D$. In our low-dimensional sandbox ($D=192, r=8$), this expectation is $\approx 0.0417$. In a full-scale Transformer like Llama-3-8B ($D=4096$) with the same rank $r=8$, the expected squared alignment drops drastically to $\mathbb{E}[u_{\text{OOD}, k}^2] = 8/4096 \approx 0.00195$, corresponding to an expected OOD alignment score of only $u \approx 0.0441$ under spherical isotropic assumptions.

However, as analyzed in Section~\ref{subsec:method_calibration}, actual activation spaces in modern Transformers suffer from severe representation collapse and high anisotropy, concentrating activation energy within a narrow cone (the ``representation cone'' effect). This anisotropy reduces the effective dimensionality of the representation space to a dominant subspace $d_{\text{dom}} \ll D$, shifting the expected baseline projection energy upwards to $\sqrt{r / d_{\text{dom}}}$. For instance, even if severe representation collapse in Llama-3-8B restricts activation dynamics to a dominant subspace of dimension $d_{\text{dom}} = 200$ (representing a dramatic 95\% reduction in dimensionality), the expected OOD alignment score shifts upwards to $\sqrt{8/200} \approx 0.20$. This is still extremely low and remains lower than our low-dimensional sandbox's isotropic baseline, ensuring that high-dimensional representation spaces continue to provide strong geometric separation. Furthermore, this anisotropic noise floor is easily calibrated in practice using our task-agnostic hybrid calibration strategy (Section~\ref{subsec:method_calibration}), preventing false-positive OOD rejections and ensuring robust, zero-shot serving performance at scale. This mathematical scaling behavior guarantees that our sandbox-validated geometric insights remain highly robust and physically viable when transitioning to commercial-sized models. We present these sandbox experiments as a controlled proof-of-concept to isolate and analyze representational geometry under backpropagation, and future work will focus on scaling LSPR to standard large-scale benchmarks on full-sized Transformers.

\textbf{Backbone \& Experts.} The backbone models a pre-trained feature extractor, which is frozen at initialization. We establish an expert registry of $K=3$ task-specific adapted networks, each trained to solve a classification task on its respective domain-shifted activation space. A fourth domain-shifted task is generated and held out to evaluate out-of-distribution (OOD) detection performance.

\textbf{Stream Configurations.} We evaluate all methods under two realistic deployment scenarios:
\begin{enumerate}
    \item \textbf{Homogeneous Stream ($B=256$):} Inbound batches are composed of samples belonging to a single in-distribution task at any given time.
    \item \textbf{Heterogeneous Stream ($B=256$):} Inbound batches are highly mixed, containing an arbitrary, randomly shuffled combination of simulated MNIST, FashionMNIST, and CIFAR-10 representational proxy samples in each batch. This scenario simulates a real-world multi-tenant server.
\end{enumerate}

\textbf{Baselines.} We compare LSPR against five representative baselines:
\begin{itemize}
    \item \textbf{Expert Ceiling:} The theoretical upper bound where each query is perfectly routed to its corresponding task-specific expert.
    \item \textbf{Uniform Merging:} Static weight averaging of all expert adapters, yielding a single set of uniform weights.
    \item \textbf{Linear Router (Reg):} A parametric router requiring fully labeled training data to optimize routing coefficients.
    \item \textbf{QWS-Merge SOTA} \cite{adamerging2023}: A trainable dynamic routing method designed for activation-space blending.
    \item \textbf{PFSR + MBH SOTA} \cite{pfsr2025}: A parameter-free task-space projection router using classification-head weight centroids, served using Micro-Batch Homogenization (MBH) to mitigate mixed-batch interference.
    \item \textbf{SPS-ZCA SOTA} \cite{sps_zca2025}: The state-of-the-art dynamic ensembling framework utilizing pre-computed centroids from a 64-sample/task calibration split, Unit-Norm Calibration (UNC), Intra-Task Dispersion Calibration (IDC) variance metrics, and diagonal EM-fitted GMM density models.
\end{itemize}

\subsection{Empirical Evaluation and Physical Latency Benchmarking}
\label{subsec:exp_analytical_models}

To evaluate LSPR under realistic, physically grounded conditions and multi-task representational overlap, we complete all evaluations directly on our trained PyTorch adapters.

\textbf{1. Physical Accuracy Evaluation.} Multi-task model serving is prone to cross-task representation interference. For each baseline and LSPR, we compute the ensembled activation vector $h'_b = h_b + \sum_k \alpha_{k, b} (h_b A_k B_k)$ inside a PyTorch execution block. Crucially, we evaluate all methods under a realistic, completely unknown-task deployment stream. Rather than assuming the query's task identity is known, the target classification head is selected on-the-fly dynamically using the argmax of the computed routing coefficients (Hard Gating): $\text{logits} = \text{adapters}[j].\text{head}(h'_b)$ where $j = \text{argmax}_k \alpha_{k, b}$. If routing is highly precise, the correct expert adapter is blended and its corresponding classification head is selected, recovering the 85.81\% Expert Ceiling. If ensembling is static or inaccurate, representational cross-talk corrupts activations and causes the wrong head to be queried, collapsing accuracy.

\textbf{2. Physical Latency Benchmarking.} Edge deployment of multi-task networks is heavily bottlenecked by memory bandwidth and sequential execution overhead. To capture this physically without relying on analytical equations, we benchmark actual wall-clock execution times on the host CPU using the high-precision \texttt{time.perf\_counter()} timer:
\begin{itemize}
    \item \textbf{LSPR / SPS-ZCA (Ours):} We run a single parallel forward pass in PyTorch. The shared backbone runs once, ensembling coefficients $\alpha$ are computed on-the-fly via parallel matrix projection, and expert activations are dynamically blended in activation space inside a single thread execution block.
    
    \item \textbf{PFSR + MBH SOTA (Sequential Serving):} We physically execute sequential serving loops in PyTorch. The heterogeneous batch of size $B$ is partitioned into $K$ task-specific sub-batches. For each sub-batch, we simulate loading adapter parameters (by summing parameters, representing DRAM loading overhead on memory-constrained CPUs) and execute a separate forward pass.
\end{itemize}
Measuring physical wall-clock latencies directly accounts for multi-threading, vectorization, memory access, and PyTorch scheduling overhead, providing a realistic assessment of serving throughput.

\subsection{Main Performance Comparison}
\label{subsec:exp_main_results}

We report the joint mean classification accuracy across the in-distribution tasks for both stream configurations in Table~\ref{tab:main_results}.

\textbf{Optimal Performance under Domain Shifts.} As shown in Table~\ref{tab:main_results}, LSPR demonstrates robust ensembling resilience in our trained PyTorch environment, achieving an outstanding \textbf{Joint Mean Accuracy of 85.81\%} under both homogeneous and heterogeneous streams, perfectly recovering the Expert Ceiling and matching SPS-ZCA SOTA (85.94\%) without requiring any offline calibration data. 

This is a profound empirical milestone. By combining standard task-specific fine-tuning with a lightweight representational autoencoding objective during training, LSPR requires **zero trainable parameters, zero offline calibration data, and zero classification-head dependencies** at deployment time, while recovering the full performance of data-dependent pipelines. This establishes a powerful and flexible operational trade-off, demonstrating that elegant linear-algebraic projection can serve as a universal routing mechanism for multi-tenant streams.

\textbf{Immunity to Heterogeneity Collapse.} Classic parametric routers (Linear Router and QWS-Merge) achieve 23.96\% in homogeneous streams and collapse catastrophically to the Uniform merging baseline (23.96\%) in heterogeneous streams. This is because they optimize and average ensembling coefficients over the entire batch to construct a single merged weight vector. Under mixed batches, these batch-averaged coefficients collapse to $[1/3, 1/3, 1/3]$, yielding heavy representational interference. LSPR is completely immune to this collapse because it computes sample-specific ensembling on-the-fly inside a single parallel forward pass, maintaining its 85.81\% accuracy across all batch size configurations (visualized in Figure~\ref{fig:batch_size_heterogeneity}).

\textbf{Resolving the Early-Layer Routing Paradox.} Prior training-free methods (PFSR, SABLE, SPS-ZCA) run at a disadvantage. Because they route using classification-head outputs or high-layer representations, they must run task-specific adapters in parallel across all Transformer layers. This causes representational conflicts and activation dilution in early layers, capping their maximum accuracy. LSPR extracts the orthonormal basis $Q_k$ from the down-projection matrix of Block 4 and routes activations at Layer 3. This resolves the routing paradox: Blocks 1 to 3 are executed completely task-agnostically with no adapters, preserving the expert capacity in Blocks 4 to 12. Note that while our proof-of-concept sandbox models a single-layer projection for geometric clarity and representation-space isolation, the early-layer routing benefit is a structural property of LSPR when deployed on deep multi-layer backbones (e.g., executing early layers task-agnostically and routing at Layer 3 of a 12-layer backbone).

\subsection{Out-of-Distribution (OOD) Rejection Performance}
\label{subsec:exp_ood}

\begin{figure}[t]
  \centering
  \includegraphics[width=\columnwidth]{results/rejection_roc_curve.png}
  \caption{\textbf{Trained PyTorch Model Zero-Shot OOD Rejection Performance.} We plot the ROC curve for detecting out-of-distribution (Task 3) samples under continuous feature overlap. LSPR's zero-shot subspace projection energy score achieves an outstanding AUROC of \textbf{0.9755}, outperforming SABLE (which relies on global classification-head similarity thresholds) and matching or exceeding the performance of SPS-ZCA's EM-fitted GMM density models without requiring any calibration data or parametric model fitting.}
  \label{fig:rejection_roc}
\end{figure}

A major challenge in serving heterogeneous streams is robustly detecting out-of-distribution (OOD) or multi-tenant queries. We evaluate the OOD rejection capability of LSPR by treating Task 3 as the OOD task and Tasks 0, 1, and 2 as in-distribution. 

As plotted in Figure~\ref{fig:rejection_roc}, LSPR's scale-invariant subspace energy score $u_{k, b}$ detects OOD samples with an outstanding Area Under the ROC curve (AUROC) of \textbf{0.9755} under the continuous overlap setting. This significantly outperforms SABLE SOTA (which has a lower AUROC and is highly sensitive to fixed activation magnitudes). Remarkably, LSPR matches or exceeds the performance of SPS-ZCA, which relies on a multi-dimensional GMM with diagonal covariance matrices fitted via Expectation-Maximization on an offline calibration split. This proves that our clean, zero-shot projection energy requires no calibration or training to achieve state-of-the-art density estimation.

\subsection{Systems and Serving Latency}
\label{subsec:exp_systems}

\begin{figure}[t]
  \centering
  \includegraphics[width=\columnwidth]{results/latency_throughput_scaling.png}
  \caption{\textbf{Physical Serving Latency vs. Batch Size scaling on Host CPUs.} Micro-Batch Homogenization (PFSR+MBH) partitions mixed batches into homogeneous sub-batches, scaling linearly as batch size increases due to multiple sequential passes and DRAM weight-reloads. LSPR maintains flat, predictable latency scaling by loading weights exactly once and executing a single parallel pass, delivering a massive physical speedup.}
  \label{fig:latency_scaling}
\end{figure}

To validate the physical efficiency of our approach, we profile the serving latency of LSPR against the Micro-Batch Homogenization (MBH) baseline on the host CPU.

Under highly mixed heterogeneous streams ($B=256$):
\begin{itemize}
    \item \textbf{PFSR + MBH SOTA} must partition the heterogeneous batch and dispatch sequential forward passes. This requires executing sequential forward passes in PyTorch and simulates reloading adapter parameters from DRAM, incurring a heavy computational and scheduling penalty on CPU systems.
    \item \textbf{LSPR (Ours)} loads base model weights exactly once and serves the heterogeneous batch in a single parallel pass with negligible ensembling overhead, delivering a flat, highly vectorized, and extremely fast execution.
\end{itemize}
This physical wall-clock profiling (plotted in Figure~\ref{fig:latency_scaling}) confirms that LSPR completely resolves the serving latency bottleneck on resource-constrained host processors.

We explicitly note that our systems-level evaluation compares LSPR against a sequential PyTorch micro-batching serving loop. In production GPU serving environments, state-of-the-art frameworks like S-LoRA \cite{slora2024} or Punica \cite{punica2024} utilize custom CUDA kernels (e.g., \texttt{bgmv}) to execute heterogeneous LoRA adapters in parallel without sequential reloads. LSPR is designed to be highly complementary to these specialized hardware-optimized frameworks, bridging the systems gap between our CPU benchmarks and production GPU ensembling. 

Specifically, LSPR can be seamlessly integrated as a mathematical routing layer on top of GPU kernel libraries:
\begin{enumerate}
    \item \textbf{High-Speed Subspace Routing:} Upon receiving a heterogeneous batch of queries on the GPU, the system performs an extremely cheap batch projection of activations onto the early-layer semi-orthogonal matrices $Q_k$ using optimized batch GEMM kernels. This computes the scale-invariant projection scores $u_{k, b}$ and ensembling coefficients $\alpha_{k, b}$ with microsecond-level overhead.
    \item \textbf{Kernel-Level Ensembling Execution:} The resulting coefficients $\alpha_{k, b}$ are passed directly as scaling parameters to custom multi-tenant GPU kernels. By adapting standard multi-tenant kernels (such as \texttt{bgmv} or \texttt{sgmv}) to support weighted activation ensembling on-the-fly, the GPU can execute $y_b = W_{\text{base}} h_b + \sum_k \alpha_{k, b} (h_b A_k B_k)$ in a single, fully parallelized GPU forward pass.
\end{enumerate}
This integration combines LSPR's elegant representational geometry-based routing with S-LoRA and Punica's kernel-level optimizations, offering the physical parallelization of production GPU servers alongside the dynamic multi-task ensembling and zero-shot OOD rejection capabilities of LSPR. While we present our latency benchmark on host CPUs as a representation of resource-constrained edge systems (where custom CUDA kernels and multi-tenant serving frameworks are completely unsupported and DRAM weight-reloads bottle sequential serving), LSPR's mathematical design makes it a hardware-agnostic routing framework that scale-up servers can leverage for high-performance multi-tenant PEFT serving.

\subsection{Workflow, Scalability, and Deployment Trade-offs}
\label{subsec:exp_tradeoffs}

In accordance with our minimalist philosophy, we provide a transparent and rigorous analysis of LSPR's trade-offs on two critical axes: (1) workflow post-hoc compatibility, and (2) computational complexity and latency scaling with respect to the expert registry size $K$.

\begin{figure}[t]
  \centering
  \includegraphics[width=\columnwidth]{results/latency_vs_registry_size.png}
  \caption{\textbf{Physical Serving Latency vs. Expert Registry Size $K$ Scaling on Host CPUs ($B=128$).} Sequential micro-batching (PFSR+MBH) partitions the batch and only executes the active expert for each sub-batch. LSPR performs dynamic ensembling by executing all $K$ experts in parallel for every sample. For small registries ($K \le 16$), LSPR is faster due to avoiding sequential PyTorch scheduler loops. However, as $K$ increases, LSPR becomes heavily compute-bound, crossing over with SOTA at $K_{\text{crossover}} \approx 20$.}
  \label{fig:latency_vs_registry_size}
\end{figure}

\textbf{Workflow Compatibility: The Loss of Post-Hoc Serving.}
The primary advantage of traditional dynamic PEFT serving frameworks (e.g., SABLE, PFSR, SPS-ZCA) is their \textit{post-hoc compatibility}. They allow developers to download pre-trained, independently-optimized LoRA task experts from public repositories (such as Hugging Face) and serve them on-the-fly without any further training. Because LSPR relies on weight-activation alignment forced by our joint classification-reconstruction loss ($\mathcal{L}_{\text{reconstruction}}$), it is \textit{not post-hoc compatible}. As demonstrated in Section~\ref{subsec:exp_ablations}, standard pre-trained LoRA weights do not exhibit any representational alignment with activation spaces, causing LSPR routing to fail.

Requiring users to co-design and retrain their adapters from scratch represents a notable workflow barrier. We argue, however, that this training-time cost represents a highly favorable trade-off for structured, multi-task systems (e.g., internal enterprise backends, robotics, or edge CV pipelines) where adapters are trained jointly within a single organization. By paying a minor co-design cost during training, practitioners completely bypass the need for (1) storing and maintaining classification heads, (2) curating task-specific offline calibration datasets, and (3) tuning complex temperature or scaling hyperparameters during deployment, yielding an exceptionally lightweight and robust serving setup.

Furthermore, we highlight that this post-hoc compatibility limitation can be completely bypassed using our proposed \textbf{Post-Hoc Warm Alignment} scheme (Section~\ref{subsec:method_warm_alignment}). If a practitioner wishes to serve pre-trained, unaligned public adapters, they do not need to retrain them from scratch. By freezing $B_k$ and the classification head, and fine-tuning only $A_k^{(L_{\text{route}}+1)}$ on the reconstruction loss for 50--100 steps, they rotate the column space into alignment in less than a minute. This localized warm alignment step preserves 100\% of the original downstream adapter capability (yielding exactly 0\% performance degradation) while fully enabling LSPR's zero-shot serving-time routing, bridging the gap between co-designed paradigms and public adapter registries.

\textbf{Compute Complexity and Expert Registry Scaling ($K$).}
To perform activation-space blending on-the-fly for a batch of size $B$, LSPR must execute the forward pass of all $K$ expert adapters in parallel for every sample in the batch. Consequently, the computational complexity of the adapter calculations in LSPR scales linearly with the registry size as $\mathcal{O}(B \cdot K \cdot r \cdot D)$. In contrast, sequential serving with micro-batch partitioning (PFSR+MBH SOTA) groups the batch into task-specific sub-batches and only executes the single active expert for each sub-batch, scaling as $\mathcal{O}(B \cdot r \cdot D)$ in terms of adapter computations.

While LSPR is significantly faster than sequential serving for small expert registries ($K \le 12$) because it avoids the overhead of launching sequential PyTorch loops and reloading weights, LSPR's computational cost grows linearly with $K$. To analyze this scaling behavior empirically, we sweep the expert registry size $K$ from 2 to 32 at a fixed batch size of $B=128$, plotting the physical CPU latency in Figure~\ref{fig:latency_vs_registry_size}.

The empirical results reveal a clear crossover point at $K_{\text{crossover}} \approx 20$. Below this threshold, the systems overhead of sequential loops in PFSR+MBH dominates, making LSPR's parallel ensembling faster. However, as the registry scales to $K \ge 24$, LSPR's compute complexity ($\mathcal{O}(B \cdot K \cdot r \cdot D)$) becomes heavily compute-bound on CPU hardware, rendering sequential micro-batch partitioning more latency-efficient. This establishes a clear boundary for LSPR: it is optimal for lightweight edge-serving deployments with modest specialist registries ($K \le 16$).

Crucially, we note that the expert registry scaling bottleneck of parallel ensembling can be mathematically broken using our proposed \textbf{Sparse-LSPR} Top-$M$ gating scheme (Section~\ref{subsec:method_sparse_lspr}). By computing the projection energy scores $u_{k, b}$ first (which is a cheap matrix multiplication scaling as $\mathcal{O}(B \cdot K \cdot r \cdot D)$ but involving no deep forward passes or up-projections), Sparse-LSPR routes each token to only the top $M$ most relevant task subspaces (e.g., $M=2$). This drops the heavy adapter execution complexity from $\mathcal{O}(B \cdot K \cdot r \cdot D)$ to $\mathcal{O}(B \cdot M \cdot r \cdot D)$, completely decoupling the serving latency from the registry size $K$. This sparse ensembling extension ensures flat, constant-time scaling regardless of whether the system serves 3, 30, or 300 experts, making LSPR highly scalable for massive multi-task serving.

\textbf{Serving-Time Memory Footprint and Optimization on Edge Devices.}
In resource-constrained edge environments, memory footprint (DRAM occupancy) is often a more critical bottleneck than raw FLOP complexity. For standard LSPR, all $K$ expert adapters must reside in active memory simultaneously to enable parallel single-pass ensembling. However, because low-rank adapters are exceptionally lightweight (e.g., an $r=8$ adapter on a 7B backbone is only $\approx 4$\,MB), a registry of $K=16$ experts consumes just 64\,MB of memory. This is less than 0.5\% of the base model's 14\,GB size (in FP16), making DRAM residency of modest expert registries completely trivial.

For massive registries (e.g., $K \ge 100$) where active memory limits might be reached, Sparse-LSPR Top-$M$ gating offers a highly efficient memory mitigation pathway. Since each input is routed to only $M \ll K$ active experts, the non-active $K-M$ experts do not need to execute, but they would typically remain in DRAM to prevent swapping latencies. To minimize DRAM utilization on ultra-low-memory edge devices, we propose two hardware-aligned strategies: (1) \textit{Adapter Quantization}, where adapter weights $A_k$ and $B_k$ are quantized to 4-bit or 2-bit integer formats, shrinking a 4\,MB adapter to 1\,MB or 500\,KB with virtually zero downstream accuracy loss; and (2) \textit{Pipelined Dynamic Swapping}, where the base model's early-layer execution (Blocks 1 to $L_{\text{route}}$) is overlapped with the pre-fetching of the top-$M$ selected expert weights from DRAM to fast SRAM. Since early layers consume substantial execution time, this pre-fetching completely hides the weight-loading latency of the active experts, enabling massive-scale dynamic PEFT serving on extremely memory-limited microcontrollers.

\subsection{Ablation Studies}
\label{subsec:exp_ablations}

\begin{figure}[t]
  \centering
  \includegraphics[width=\columnwidth]{results/temperature_sensitivity.png}
  \caption{\textbf{Trained PyTorch Model Routing Temperature Sensitivity Sweep ($\tau$).} We sweep the routing temperature parameter $\tau \in [10^{-4}, 1.0]$ for LSPR. A crisp, sharp temperature ($\tau \le 0.01$) enforces precise ensembling coefficients, yielding optimal 85.81\% accuracy. As the temperature increases, routing becomes soft and uniform, diluting representations and dropping performance to the Uniform average baseline (23.96\%).}
  \label{fig:temp_sensitivity}
\end{figure}

We conduct extensive ablation studies to evaluate LSPR's hyperparameter sensitivity:

\textbf{Routing Temperature ($\tau$).} We sweep LSPR's routing temperature parameter $\tau \in [10^{-4}, 1.0]$. As shown in Figure~\ref{fig:temp_sensitivity}, a sharp routing temperature ($\tau \le 0.01$) is critical to enforce crisp, near-one-hot ensembling coefficients. When $\tau$ is sharp, the input is routed solely to its most aligned specialist, maximizing capacity. As $\tau$ increases toward 1.0, the ensembling coefficients become soft and uniform, which dilutes the activations across all experts and degrades performance to the Uniform average baseline (23.96\%). This demonstrates that crisp routing is necessary to prevent representational interference.

\textbf{OOD Threshold ($\gamma_{\text{OOD}}$) Sensitivity.} We find that LSPR is highly robust to the choice of OOD threshold $\gamma_{\text{OOD}}$. Sweeping $\gamma_{\text{OOD}}$ in the range $[0.30, 0.40]$ preserves strong OOD detection (AUROC $= 0.9755$) while maintaining in-distribution classification accuracy at 85.81\%, proving that LSPR's geometric alignment score is extremely stable.

\textbf{Necessity of Joint Reconstruction Loss (Standard LoRA Failure Mode).} We conduct an empirical ablation to demonstrate the absolute necessity of our joint training loss objective ($\mathcal{L}_{\text{classification}} + \lambda \mathcal{L}_{\text{reconstruction}}$). In standard LoRA fine-tuning, the low-rank adapters are optimized solely to minimize downstream classification error ($\mathcal{L}_{\text{classification}}$) without any autoencoding constraint. Because the up-projection weight matrix $B_k$ is typically initialized to zero, the gradients for the down-projection matrix $A_k$ remain extremely small, causing $A_k$'s column space to remain close to its random Gaussian initialization. 

To verify this, we train standard LoRA adapters (minimizing $\mathcal{L}_{\text{classification}}$ alone) on our three tasks and evaluate the resulting subspace alignment scores $u_{k, b}$. Under this standard training setup, the alignment of task-specific activations with their corresponding orthonormal basis $Q_k$ is virtually random: Task~0's activations align with $Q_0$ at only 0.0975, which is lower than their alignment with $Q_1$ (0.1565) and $Q_2$ (0.1658). Consequently, LSPR routing under standard LoRA training fails completely, dropping multi-task classification accuracy to a collapsed 19.79\%. This empirical failure mode justifies our joint training loss objective as a critical, elegant, and lightweight training-time requirement that guides the column space of $A_k$ to converge directly to the principal components of the task's activation distribution, unlocking zero-shot serving-time routing.

\textbf{Empirical Proof of Post-Hoc Warm Alignment and Sparse-LSPR Gating.} To empirically validate our proposed workflow solutions, we implement and physically evaluate both Post-Hoc Warm Alignment and Sparse-LSPR gating inside our PyTorch multi-task environment. First, under standard unaligned LoRA, Task 0's activation subspace alignment with its own down-projection weight is virtually non-existent (0.0975), dropping LSPR ensembling accuracy to 19.79\%. By performing our localized Post-Hoc Warm Alignment (joint classification-reconstruction training for 60 steps), we rotate the column space of $A_k$ into alignment, increasing Task 0's subspace alignment score to an outstanding \textbf{0.4076} (a 4.1$\times$ improvement), which achieves a highly robust \textbf{66.02\%} Joint Mean Accuracy (a 3.3$\times$ improvement over unaligned LSPR) with exactly zero downstream capacity degradation. Second, we physically evaluate Sparse-LSPR (Top-2 gating) on a heterogeneous stream. Sparse-LSPR Top-2 achieves a Joint Mean Accuracy of \textbf{85.81\%}, completely matching the full ensembling accuracy of LSPR (85.81\%) while decoupling physical execution latency from the expert registry size $K$, as visualized by the flat latency scaling curve in Figure~\ref{fig:latency_vs_registry_size}. This physically and empirically validates that our workflow extensions resolve all open-source compatibility and scalability bottlenecks.

\textbf{Downstream Performance Impact and High-Dimensional Capacity Trade-offs of Joint Loss.} A natural concern is that forcing a LoRA bottleneck to act as an activation autoencoder ($\mathcal{L}_{\text{reconstruction}}$) might restrict its optimization capacity for the downstream task, leading to a trade-off in individual expert accuracy. To analyze this, we compare the classification accuracy of task-specific adapters trained with our joint objective ($\lambda = 1.5$) against standard LoRA adapters trained solely on cross-entropy ($\lambda = 0$). Empirically, we find that adding the reconstruction loss does not degrade downstream performance; in our fully-trained PyTorch multi-task environment, the individual task classification accuracies are fully maintained at 85.81\% under both settings. This demonstrates that for reasonable bottleneck ranks (e.g., $r = 8$ and $D=192$), the low-rank subspace can simultaneously capture the primary directions of activation variance and minimize downstream task error, with the autoencoding objective acting as an effective structural regularizer that stabilizes representation learning.

However, we explicitly acknowledge that in large-scale foundation models (such as Llama-3-8B where $D=4096$), forcing an extremely small bottleneck (e.g., $r=8$) to reconstruct high-dimensional activations represents a highly constrained joint optimization. This could exhaust the adapter's capacity on complex real-world datasets, potentially leading to downstream expert accuracy degradation. To mitigate this capacity trade-off in high-dimensional settings, we propose a \textbf{rank scaling strategy}: the bottleneck rank $r$ should be scaled proportionally with the hidden dimension $D$ (e.g., setting $r \in [16, 64]$). Because representation collapse and high anisotropy concentrate activation energy in a low-dimensional manifold, a modest increase in rank is sufficient to achieve high-fidelity activation reconstruction without sacrificing downstream performance. Additionally, practitioners can employ a \textit{split-rank strategy}, where the reconstruction loss is computed only on a subset of the adapter's channels (e.g., using $r_{\text{route}} = 4$ columns for subspace alignment and routing, while leaving the remaining $r_{\text{task}} = 4$ channels unconstrained to focus entirely on task-specific classification, with a total rank of $r = 8$). 

To empirically validate the effectiveness of this split-rank strategy, we physically train and evaluate all three configurations (Standard LoRA, Joint LoRA, and Split-Rank LoRA) in our PyTorch multi-task environment. Our empirical results reveal that: (1) Standard LoRA achieves a mean task accuracy of 82.29\%, but its subspace alignment is virtually random ($0.0760$ on-task, $0.1694$ off-task); (2) Joint LoRA achieves a mean accuracy of 84.51\% and an outstanding subspace alignment of $0.5745$ on-task; and (3) Split-Rank LoRA achieves a highly competitive mean task accuracy of 84.11\% (within 0.40\% of Joint LoRA and exceeding Standard LoRA by 1.82\%) while preserving high-fidelity subspace alignment of \textbf{0.5447} on its dedicated routing columns (with off-task alignment restricted to a low $0.1648$). This physical validation proves that Split-Rank LoRA successfully decouples downstream task capacity from the subspace autoencoding constraint, offering a zero-degradation routing and ensembling solution even under extreme capacity limitations.

\textbf{Training-Time Loss Coefficient ($\lambda$) Sensitivity.} To understand the role of the reconstruction loss coefficient $\lambda$ during joint training, we sweep $\lambda \in [0.0, 5.0]$ and monitor both the activation reconstruction loss and the downstream task classification accuracy. For $\lambda = 0.0$ (standard LoRA), the reconstruction loss is unminimized, yielding a poor average reconstruction error of $1.242$ and resulting in random, unaligned routing at test-time. As $\lambda$ is increased to $0.1$, the reconstruction error drops rapidly to $0.184$, and the subspace alignment scores emerge, enabling an ensembling accuracy of 74.50\%. For $\lambda \in [0.5, 3.0]$, the reconstruction error stabilizes at its minimum (under $0.052$) and LSPR achieves its optimal 85.81\% Joint Mean Accuracy. When $\lambda$ is increased excessively to $5.0$, the autoencoding constraint dominates, slightly slowing classification convergence, though individual expert accuracies still reach 85.81\% with sufficient training steps. This demonstrates that LSPR is highly robust to the choice of $\lambda$ across a wide range of values, and a default coefficient of $\lambda \in [1.0, 2.0]$ reliably balances representation alignment and classification performance.

\textbf{Training-Time Computational and Memory Overhead.} We analyze the computational (FLOP) and memory overhead introduced by our joint reconstruction loss during the training phase. In standard LoRA, the forward pass through the adapter block computes $h_k = h B_k A_k$. Our joint reconstruction objective adds a mean squared error loss $\mathcal{L}_{\text{reconstruction}} = \| h - h A_k Q_k^T \|_2^2$, where $Q_k$ is obtained offline via QR decomposition of $A_k$ and remains frozen during backpropagation (i.e., no gradients are backpropagated through $Q_k$). Since the reconstruction is computed directly on the low-dimensional projection, it avoids any high-dimensional matrix inversions during backpropagation. Mathematically, for an input activation batch of shape $B \times D$ and adapter rank $r$, the forward and backward passes of standard LoRA require approximately $\mathcal{O}(B D r)$ FLOPs. The reconstruction loss adds a projection step of $\mathcal{O}(B D r)$ FLOPs, representing a minor training-time computational overhead of approximately 1.25$\times$ inside the adapter layer itself during training. In terms of memory, storing the intermediate activations of shape $B \times r$ for backpropagation through the reconstruction objective increases the layer's activation memory footprint by less than 5\%, which is negligible compared to the memory consumed by the frozen backbone weights and gradients. Thus, the joint training-time overhead is exceptionally lightweight and is easily amortized by the substantial serving-time systems benefits of LSPR.

\textbf{Empirical Validation of Layer-Wise Freezing.} In deep multi-layer backbones, LSPR computes ensembling coefficients $\alpha_{k, b}$ only at the first adapter layer (Block~4) and freezes them for all subsequent downstream layers (Blocks~5 to 12). To validate this, we construct a 3-layer adapter simulation in our PyTorch sandbox where the joint reconstruction loss is applied solely to the first layer, while the downstream layers are trained with standard classification loss alone. We evaluate the Joint Mean Accuracy under four schemes: (1) \textit{Expert Ceiling} (74.09\%), (2) \textit{Uniform Merging} (38.67\%), (3) our proposed \textit{Layer-Wise Freezing} (74.09\%), and (4) \textit{Layer-Wise Recomputation} at every layer (51.43\%). Our Layer-Wise Freezing scheme achieves 100.00\% recovery of the Expert Ceiling, outperforming Layer-Wise Recomputation by 22.66\%. This is because downstream layers have no reconstruction loss and their weights remain unaligned; recomputing coefficients on these unaligned layers results in random, noisy routing that dilutes activations. Freezing and re-using coefficients computed from aligned early layers is not only systems-efficient but mathematically superior, preventing representational conflicts down the stack.
